\section{Quotient-Wronskian reduction}\label{sec:quotient}

Fix \(y_1<y_2<y_3<y_4\), put \(u_j(t)=\Phi(t-y_j)\), and write
\(p_j=t-y_j\), so \(p_4<p_3<p_2<p_1\). Define
\[
 v_j=(u_j/u_1)',\qquad w_j=(v_j/v_2)'.
\]
These quotients are introduced only after their denominators are known to be
positive: \(u_1>0\) by \secref{sec:kernel}, and \(v_2>0\) by the next
display. Later division by \(w_3\) is likewise made only after
\eqref{eq:gj-Lambda} proves \(w_3>0\).
Strict log-concavity gives
\[
 v_j=\frac{u_j}{u_1}A(p_j,p_1)>0\qquad(j\geq2).
\]
The case \(k=2\) of Lemma~\ref{lem:quotient-integral} below therefore proves
strict order two directly.
At the fixed orders used in this paper, successive column division and
differentiation give the exact three-stage cascade
\begin{align}
 W(u_1,u_2,u_3)&=u_1^3W(v_2,v_3),\label{eq:W3-quotient}\\
 W(u_1,u_2,u_3,u_4)&=u_1^4W(v_2,v_3,v_4),\label{eq:W4-first}\\
 W(v_2,v_3,v_4)&=v_2^3W(w_3,w_4),\label{eq:W4-second}\\
 W(w_3,w_4)&=w_3^2\frac{\dd}{\dd t}
 \left(\frac{w_4}{w_3}\right).\label{eq:W4-third}
\end{align}
Consequently
\begin{equation}\label{eq:W4-quotient}
 W(u_1,u_2,u_3,u_4)=u_1^4v_2^3w_3^2
 \frac{\dd}{\dd t}\left(\frac{w_4}{w_3}\right).
\end{equation}
No sign premise is used in these algebraic identities beyond the stated
nonvanishing needed to form the quotients.

Let \(\T\) denote simultaneous translation of all point arguments. From
\eqref{eq:AM},
\begin{equation}\label{eq:TlogA}
 \T A(a,b)=q(b)-q(a),\qquad \T\log A(a,b)=M(a,b).
\end{equation}
The sign convention in \eqref{eq:TlogA} will be used below. Since
\[
 \frac{v_j}{v_2}=\frac{u_j}{u_2}
 \frac{A(p_j,p_1)}{A(p_2,p_1)},
\]
direct differentiation gives
\begin{equation}\label{eq:gj-expanded}
 g_j:=\T\log(v_j/v_2)
 =A(p_j,p_2)+M(p_j,p_1)-M(p_2,p_1).
\end{equation}
The exact weighted-mean split
\[
 A(p_j,p_1)M(p_j,p_1)
 =A(p_j,p_2)M(p_j,p_2)+A(p_2,p_1)M(p_2,p_1)
\]
turns \eqref{eq:gj-expanded} into
\begin{equation}\label{eq:gj-Lambda}
 g_j=\frac{A(p_j,p_2)}{A(p_j,p_1)}
 \Lambda(p_j;p_2,p_1)>0.
\end{equation}
Thus \(w_j=(v_j/v_2)g_j>0\), and
\eqref{eq:W3-quotient} has the strict order-three sign.

Define
\[
 \Psi(\xi;m,r)=\Lambda(\xi;m,r)+\T\log\Lambda(\xi;m,r).
\]
Because \(w_j=(v_j/v_2)g_j\),
\begin{equation}\label{eq:L4-first}
 \T\log(w_4/w_3)
 =(g_4+\T\log g_4)-(g_3+\T\log g_3).
\end{equation}
Equation \eqref{eq:gj-Lambda} and \eqref{eq:TlogA} give
\[
 \T\log g_j=M(p_j,p_2)-M(p_j,p_1)+\T\log\Lambda_j,
\]
where \(\Lambda_j=\Lambda(p_j;p_2,p_1)\). A second use of the same mean split
gives
\[
 g_j+M(p_j,p_2)-M(p_j,p_1)=\Lambda_j-A(p_2,p_1).
\]
The common last term cancels in \eqref{eq:L4-first}, yielding the exact
identity
\begin{equation}\label{eq:L4-Psi}
 \T\log(w_4/w_3)
 =\Psi(p_4;p_2,p_1)-\Psi(p_3;p_2,p_1).
\end{equation}

\begin{lemma}[Iterated quotient integral]\label{lem:quotient-integral}
Let \(f_1,\ldots,f_k\) be continuously differentiable on an interval,
let \(f_1>0\), and put \(G_j=f_j/f_1\). For
\(t_1<\cdots<t_k\),
\begin{equation}\label{eq:iterated-determinant}
 \det[f_j(t_i)]_{i,j=1}^k
 =\prod_{i=1}^k f_1(t_i)
 \int_{t_1}^{t_2}\!\cdots\!\int_{t_{k-1}}^{t_k}
 \det[G'_{j+1}(\tau_i)]_{i,j=1}^{k-1}
 \dd\tau_1\cdots\dd\tau_{k-1}.
\end{equation}
The identity may be applied repeatedly whenever the first member of the new
derivative family is positive.
\end{lemma}

\begin{proof}
Factor \(f_1(t_i)\) from row \(i\). The first column is then one. Replace
successive rows, from bottom to top, by their difference with the preceding
row. Each remaining entry is
\(G_j(t_{i+1})-G_j(t_i)=\int_{t_i}^{t_{i+1}}G_j'(\tau)\dd\tau\).
Determinant multilinearity gives \eqref{eq:iterated-determinant}. Since
\(\tau_i\in[t_i,t_{i+1}]\), the integration nodes remain ordered, so the same
argument applies to the derivative determinant.
\end{proof}

\begin{sloppypar}
Only the instances of Lemma~\ref{lem:quotient-integral} of sizes two through
four enter the theorem. The maintained formal engine proves exactly those
instances and the corresponding strict integral boxes; no arbitrary-order
quotient theorem is a premise of Theorem~\ref{thm:main}. In particular,
\path{PF4.FixedOrderQuotientWronskian.fixedOrderFour_quotientWronskian_package}
packages \eqref{eq:W3-quotient}--\eqref{eq:W4-third} together with the reversed
endpoint orientation \(p_4<p_3<p_2<p_1\). Its independent replay boundary is
\texttt{CERT29}.
\end{sloppypar}

\begin{lemma}[Confluent divided-difference limit]\label{lem:confluent-limit}
Let \(f\in C^{2k-2}\), let
\(V(a)=\prod_{0\leq i<j<k}(a_j-a_i)\), and let \(a,b\) be fixed strictly
increasing \(k\)-tuples. Then
\begin{equation}\label{eq:two-sided-confluence}
 \lim_{h\downarrow0}
 \frac{\det[f(z+h(a_i-b_j))]_{i,j=0}^{k-1}}
 {h^{k(k-1)}V(a)V(b)}
 =\frac{(-1)^{k(k-1)/2}}
 {\prod_{r=0}^{k-1}(r!)^2}
 \det[f^{(i+j)}(z)]_{i,j=0}^{k-1}.
\end{equation}
If only the row nodes coalesce while \(y_0<\cdots<y_{k-1}\) remain fixed,
then
\begin{equation}\label{eq:one-sided-confluence}
 \lim_{h\downarrow0}
 \frac{\det[f(z+ha_i-y_j)]_{i,j=0}^{k-1}}
 {h^{k(k-1)/2}V(a)}
 =\frac{\det[f^{(i)}(z-y_j)]_{i,j=0}^{k-1}}
 {\prod_{r=0}^{k-1}r!}.
\end{equation}
\end{lemma}

\begin{proof}
Apply successive Newton divided-difference row operations, whose divisors
are the positive numbers \(h(a_j-a_i)\), and let \(h\downarrow0\). This gives
\eqref{eq:one-sided-confluence}. Applying the same operation to the columns
gives derivatives with respect to \(-y\), hence the factor
\((-1)^{0+\cdots+(k-1)}\), and proves
\eqref{eq:two-sided-confluence}. This records the factorial normalization and
does not assume a nonzero leading coefficient.
\end{proof}

\begin{proposition}[Three-point equivalence]\label{prop:equivalence}
Under the strict lower-order positivity already proved, global \(\PF_4\) is
equivalent to
\[
 \partial_\xi\Psi(\xi;m,r)\leq0\qquad(\xi<m<r).
\]
Strict inequality implies strict global \(\PF_4\).
\end{proposition}

\begin{proof}
If \(\partial_\xi\Psi\leq0\), then \(p_4<p_3\) makes the right side of
\eqref{eq:L4-Psi} nonnegative. All prefactors in
\eqref{eq:W4-quotient} are positive. The Wronskian is therefore
nonnegative, and is positive in the strict case.

For clarity, applying Lemma~\ref{lem:quotient-integral} three times to the
order-four minor gives, with \(I_i=[t_i,t_{i+1}]\),
\begin{align*}
\det[u_j(t_i)]_{i,j=1}^4
={}&\prod_{i=1}^4u_1(t_i)
\int_{I_1\times I_2\times I_3}\prod_{i=1}^3v_2(\tau_i)\\
&\quad\times\int_{\tau_1}^{\tau_2}\int_{\tau_2}^{\tau_3}
w_3(\sigma_1)w_3(\sigma_2)
\int_{\sigma_1}^{\sigma_2}
\left(\frac{w_4}{w_3}\right)'(\rho)\dd\rho
\dd\sigma_2\dd\sigma_1\dd\boldsymbol\tau.
\end{align*}
Every factor is positive in the strict case, and every integration interval
has positive length. This proves positivity of every ordered collocation minor,
not only of the confluent Wronskian.

Conversely, assume global \(\PF_4\). Fix arbitrary
\(\xi_1<\xi_2<m<r\), choose any \(t\), and choose the columns so that
\((p_4,p_3,p_2,p_1)=(\xi_1,\xi_2,m,r)\) at \(t\); explicitly,
\(y_j=t-p_j\), which preserves \(y_1<\cdots<y_4\). Coalesce arbitrary
strictly ordered row nodes at \(t\). Equation
\eqref{eq:one-sided-confluence} and global \(\PF_4\) make the resulting
translate Wronskian nonnegative. Equations \eqref{eq:W4-quotient} and
\eqref{eq:L4-Psi} therefore give
\[
 \Psi(\xi_1;m,r)\geq\Psi(\xi_2;m,r).
\]
Because both the pair \(\xi_1<\xi_2\) and the pair \(m<r\) were arbitrary,
\(\Psi(\cdot;m,r)\) is nonincreasing for every admissible \(m,r\). Smoothness
makes this equivalent to the stated differential inequality.
\end{proof}
